\section{The limit of the parameter refinement at
\texorpdfstring{$\kappa=3/4$}{kappa=3/4}}\label{gex:limit-section}

With the asymmetry fixed at $b=1/8$ and the plain-moment exponent
fixed at $\widehat\kappa=3/4$, the
perturbation is close to the limit of this particular balanced
witness estimate. At $y=0$ the zero-margin quadratic is
\[
 Q_\varpi(\delta)=(2448+6048\varpi)\delta^2+
              (-1896+11520\varpi)\delta+370-22200\varpi.
\]
Its discriminant test gives
\[
 4AC-B^2=-576(1162800\varpi^2+286020\varpi-49).
\]
The vertex remains inside $0\le\delta\le5/6$ near the present
perturbation. Therefore, at $\widehat\kappa=3/4$, this estimate cannot
support a larger perturbation past
\[
 \varpi_*=-\frac{1589}{12920}+\frac{7\sqrt{465053}}{38760}
       =0.0001711975385\ldots,
\]
whose loss-free crossing $7/8-\varpi_*/4$ is approximately
$0.8749572006$. This is the limit of the retained endpoint inequality
at $\widehat\kappa=3/4$ only. It is not a lower bound for the true
physical source, for the same method with other plain-moment
exponents, or for other arithmetic methods.

The adaptive argument in Section~\ref{gex:adaptive-section} passes this
fixed-$(b,\kappa)$ ceiling by using the previously established boundary
in the prime-slot capacity. Of its gain of about $1.18\cdot10^{-6}$
over $\theta_0$, about $1.13\cdot10^{-6}$ is already available at the
parameter limit above. The finite feedback approaches the loss-free
value $0.87495715035054$; that limit is not asserted in
Theorem~\ref{gex:main}. Further decimal tuning within this estimate
would give little additional improvement. A substantial advance needs
a stronger arithmetic estimate.

\subsection*{Supplementary verification}

The supplementary archive (repository commit to be pinned on release)
builds on OpenAI's Lean formalization of \cite{OAI26}, at the commit
pinned in the bibliography; its definitions and mathematical input
modules are unchanged. Its final theorem is the standard Mathlib
statement \mbox{\texttt{riemannZeta s $\ne$ 0}} for
$20999/24000<\Rea s$. This statement passed exact comparison against
the target and replay in the Lean kernel, with axioms exactly
\texttt{propext}, \texttt{Classical.choice} and \texttt{Quot.sound}.
The corresponding Hecke and Dirichlet declarations compiled but were
not comparator-checked. The archive also contains exact rational checks
of the first-stage endpoint certificate. It verifies $\theta_0$, not
the adaptive boundary $\theta_*$. The latter has the conventional
proof above and the recursive proof in \cite{GaoAdaptive26}, without
a new Lean, Comparator or kernel replay. The archive accompanies
and does not replace these proofs.

\subsection*{Attribution and preparation}

The cubic-theta construction, the unchanged arithmetic inputs and the
quadratic endpoint-certificate technique \cite[Lemma 20.2]{OAI26} are
due to OpenAI \cite{OAI26}, as is the Lean formalization on which the
supplementary archive builds. The contribution here is the
modified-scale low estimate, which extends \cite[Lemma 15.1 and
Proposition 15.3]{OAI26} to the perturbed geometry; the re-run of the
endpoint certificate with the perturbation retained; the extension of
the principal Euler region below $7/8$; and the recovery at the new
boundary; the extended plain-moment range and its finite adaptive
reuse, proved in \cite{GaoAdaptive26} and
Section~\ref{gex:adaptive-section}. AI tools assisted with the derivations, algebraic checks and
drafting.

\paragraph{Contemporaneous refinements.}
Liu~\cite{LIU26} and ProofCouncil~\cite{PC26} report smaller boundaries.
ProofCouncil also extends the plain moment to $\kappa\ge13/18$ and
jointly optimizes the asymmetry $b$, which is fixed at $1/8$ here.
The adaptive extension and witness balancing therefore overlap;
their gains cannot be added to those of this paper. This paper records
a conditional derivation and its exact certificates. It claims neither
the smallest known boundary nor priority for the shared extension.

\appendix
\section{Hypothesis audit}\label{gex:audit}

For each statement of \cite{OAI26} used here, the table records its
printed hypotheses, the values used in this paper, and why its proof
transfers. Part~II of \cite{OAI26} works under the supposition
\cite[(12.1)]{OAI26}, $\Delta=\beta_*-7/8\in(0,1/24]$ and
$\kappa=2\beta_*-1$, with the geometry \cite[(12.4)]{OAI26},
$\ell=1/6$. For the first stage, by contrast, $\theta_0<\beta_*\le7/8$, with
$\sigma_0=\theta$, $\widehat\kappa=3/4$ and the geometry
\eqref{gex:geometry}, $\ell=1/6+\varpi$. The second Euler region
$D_2$ of \cite[Section 10.2]{OAI26} requires $\Rea s\ge7/8$; it is
replaced by the box \eqref{gex:second-euler-box}.
The table concerns that first stage. The adaptive stage additionally
uses Lemma~\ref{gex:adaptive-plain} and its complete proof in
\cite{GaoAdaptive26}; Section~\ref{gex:adaptive-section} records the
changed capacity, endpoint and recovery payments.

{\footnotesize
\begin{longtable}{@{}>{\raggedright\arraybackslash}p{0.17\linewidth}
                   >{\raggedright\arraybackslash}p{0.25\linewidth}
                   >{\raggedright\arraybackslash}p{0.21\linewidth}
                   >{\raggedright\arraybackslash}p{0.29\linewidth}@{}}
\toprule
Statement & Printed hypotheses & Used here & Why the proof transfers\\
\midrule
\endfirsthead
\toprule
Statement & Printed hypotheses & Used here & Why the proof transfers\\
\midrule
\endhead
\bottomrule
\endfoot
Thm.\ 1.1
 & None (unconditional)
 & $\beta_*\le7/8$, as a hypothesis
 & Used as stated; this is the main reason the theorem is conditional.\\
\addlinespace[3pt]
Prop.\ 2.1
 & $\sigma_0\in(1/2,1)$, $\Delta_0=\beta_*-\sigma_0>0$;
   common $0<\omega<\Delta_0$, $\sigma>0$;
   $\sup_{\Rea s>\sigma_0}|H_\eta-1|\le1/2$ (2.2); (2.4), (2.5)
 & $\sigma_0=\theta$, $\omega=\varepsilon_{\rm low}=\Gamma/2$;
   $H_\eta$ of \eqref{gex:principal-nonvanishing}
 & The proof uses (2.2) only for boundedness and nonvanishing of
   $H_\eta$; Section~\ref{gex:recovery-section} supplies both on every
   closed sub-half-plane of $\Rea s>2/3$.\\
\addlinespace[3pt]
Prop.\ 11.3
 & $\sigma_0\in[1/2,1)$; primitive Hecke $L$ zero-free on
   $\Rea s>\sigma_0$
 & $\sigma_0=\theta$
 & Applied as stated.\\
\addlinespace[3pt]
Lemmas 4.5, 4.6, 4.10
 & Smooth profiles; fixed $\sigma_0>0$ in Lemma 4.10
 & As stated
 & No zero-free or geometry hypothesis.\\
\addlinespace[3pt]
Lemma 4.9
 & $a\in[1/2,1]$, $0<e<10^{-3}$, zero-free disk; global form for
   $b\ge\beta_*$, $v>0$
 & $a\ge51/100$; $b=\beta_*$ or $b=7/8+8e$
 & Applied as stated; the global form needs only $b\ge\beta_*$.\\
\addlinespace[3pt]
Lemma 7.1, (7.11)--(7.12); (16.3), (16.8)
 & Regions (7.14) and (7.15), the latter with $\Rea s\ge7/8$
 & Box \eqref{gex:second-euler-box}, $\Rea s\ge\theta$
 & The local defects remain summable there
   by \eqref{gex:local-euler-defect} and
   \eqref{gex:principal-selected-factor}.\\
\addlinespace[3pt]
Lemmas 10.3, 10.4
 & Def.\ 10.1 with $D_1(\epsilon_0)$, $D_2$; $\sigma_0\in[7/8,1)$,
   $\beta_*>\sigma_0$; $z_0=17/50$
 & $\sigma_0=\theta$, $\theta<\beta_*\le7/8$;
   geometry \eqref{gex:geometry}
 & Central contours lie in $D_1(10e)$, which does not involve $7/8$;
   $\sigma_0$ enters only as the reference $C(\sigma_0)$ in (10.15);
   the bin ceiling $a\le\beta_*$ is unchanged.\\
\addlinespace[3pt]
Lemma 10.5
 & $\sigma_0\in[7/8,1)$, $\beta_*>\sigma_0$; (10.17) on a box in $D_2$;
   (10.18)
 & $\sigma_0=\theta$; the same box, now inside
   \eqref{gex:second-euler-box}
 & $D_2$ enters only through holomorphy and (10.17), proved on the
   enlarged box; (10.18) holds with $\mu=\kappa_P$ by
   \eqref{gex:principal-selected-factor}.\\
\addlinespace[3pt]
Lemma 10.6
 & $\sigma_0\in[7/8,1)$, $\beta_*>\sigma_0$; (10.20) on
   $(\beta_*+e,1/2,17/50)$
 & Small rows on $(7/8+8e,1/2,17/50)$; large rows on $(2,2,z_\infty)$
 & The line is right of $\beta_*$ by Thm.\ 1.1 and lies in $D_1(3/8)$;
   its extra cost $7/8-\beta_*+7e$ is paid from $m_{\rm small}$.\\
\addlinespace[3pt]
Lemma 11.1
 & $\sigma_0\in(1/2,1)$, $\Delta_0>0$; common $m>0$, $\omega<\Delta_0$
 & $\sigma_0=\theta$
 & Applied as stated.\\
\addlinespace[3pt]
Lemma 14.3; Prop.\ 15.2
 & Bounded log-lengths; $Q,Y'\ge1$, $P_a=Y'^2/Q\ge1$
 & Perturbed scales
 & Hypotheses verified by \eqref{gex:uniform-low-scale-margins}.\\
\addlinespace[3pt]
Lemma 15.1
 & Geometry (12.4), $\ell=1/6$
 & $0\le\ell\le1/5$
 & Proof repeated in Lemma~\ref{gex:completed-low}; $\ell$ enters only
   through $N_*=1+\ell'$ and $z_a\le\ell'$.\\
\addlinespace[3pt]
Prop.\ 15.3
 & Geometry (12.4), $\ell=1/6$; exponent $3/16$
 & Geometry \eqref{gex:geometry}
 & Proof repeated in Proposition~\ref{gex:physical-low}; the extra
   term $(5\ell-1+d)_+/8$ is paid by the tuple count, since
   $\ell<1/5$.\\
\addlinespace[3pt]
Lemma 16.2
 & First lines $(\beta_*+e,1/2,17/50)$ with $\beta_*>7/8$
 & $(7/8+8e,1/2,17/50)$
 & The proof uses only $\Rea s\ge7/8+e$ and $D_1(3/8)$.\\
\addlinespace[3pt]
Prop.\ 8.3; Lemmas 17.1, 17.6
 & Detector bins; bounded ranges; strict moment inequalities
 & As stated, with slots from $\ell=1/6+\varpi$
 & No zero-free hypothesis; slot lengths enter only through bounded
   ranges and the mesh.\\
\addlinespace[3pt]
Lemma 18.1
 & $3/4\le\kappa\le1$; if $z>0$ and $\kappa<1$, then
   $\beta_*\le(1+\kappa)/2$
 & $\kappa=3/4$; $\beta_*\le7/8$ by Thm.\ 1.1
 & Applied as stated.\\
\addlinespace[3pt]
Prop.\ 19.2
 & Under (12.1): $\Delta>0$, $\kappa=2\beta_*-1$; $\alpha=5/6$;
   prime supply above $7/37$; loss $+\Delta/4$
 & $\widehat\kappa=3/4$; supply $\ell/d>1/5$
 & The proof uses $\kappa$ only through $\delta\le\kappa$ and
   Lemma 18.1, both supplied by Thm.\ 1.1; the loss came from
   $6\kappa=9/2+12\Delta$ and vanishes.\\
\addlinespace[3pt]
Lemma 20.1
 & Lemma 10.4 with $\sigma_0=7/8$; (12.2); geometry (12.4)
 & $\sigma_0=\theta$; exponent \eqref{gex:high-central-exponent}
 & Only the affine exponent (20.4) changes; the contour and
   conductor-allocation steps use the slot system, not $\ell$.\\
\addlinespace[3pt]
Lemma 20.2
 & Geometry (12.4), $\varpi=0$; $0\le\delta\le5/6$, $0\le x\le1/2$
 & $\varpi=1/6000$
 & Same algebra with $\varpi$ retained; Lemma~\ref{gex:endpoint}
   gives margin $1/200000$.\\
\end{longtable}
}

\paragraph{Adaptive-stage hypotheses.}
At the first adaptive step, the input boundary is $\theta_0$ and
$\kappa_0=2\theta_0-1$; at the second it is the established $\theta_1$
and $\kappa_1=2\theta_1-1$. Both lie in $[13/18,3/4]$.
The global bound $\beta_*\le(1+\kappa_j)/2$ comes from the previous
stage, including its all-height and full-family continuation.
The common masks, character zeros, finite slot coefficients and
moving-radical width remain those of the extended moment.
The small-row contour stays at $7/8+8e$, and its additional shift
cost is paid explicitly. Thus no use of the output boundary is hidden
in a prime estimate, recursive child or contour choice.

\begin{thebibliography}{9}
\bibitem{GaoAdaptive26}
Z.~Gao,
\emph{The adaptive plain-moment estimate: supplementary proof},
supplementary material to this paper, 9 October 2026.
File: \nolinkurl{seven_eighths_adaptive_moment.pdf}.
\bibitem{OAI26}
\raggedright
OpenAI,
\emph{The Quasi-Riemann Hypothesis: A Zero-Free Half-Plane
$\Rea s>7/8$}, unrefereed preprint, 30 September 2026.
Version used: \texttt{adc7f1241b42e322a6451854ab7e4b4c146bf78a}.
\url{https://github.com/openai/math/blob/adc7f1241b42e322a6451854ab7e4b4c146bf78a/preprints/The-Quasi-Riemann-Hypothesis-September-30-2026/paper.pdf}.
\bibitem{DAV00}
H.~Davenport,
\emph{Multiplicative Number Theory}, 3rd ed., revised by
H.~L.~Montgomery, Graduate Texts in Mathematics~74, Springer,
New York, 2000.
\bibitem{KIN26}
\raggedright
Shiva Kintali,
\emph{A Short Proof of the Quasi-Riemann Hypothesis},
preprint, revised 27-page version served at the URL below on
8 October 2026, SHA-256
\texttt{f7ddc46632e37466e85951ea67acbd3c}\allowbreak
\texttt{2abc19c8259fac51be569b1f6bbbe8e5};
equation numbers refer to this version.
\url{https://shivakintali.github.io/papers/QRH.pdf}.
\bibitem{LIU26}
\raggedright
B.~Liu,
\emph{Slightly improved zero-free half-planes for the quasi-Riemann
hypothesis}, preprint, arXiv:2610.12234v1, 8 October 2026.
\url{https://arxiv.org/abs/2610.12234v1}.
\bibitem{PC26}
\raggedright
ProofCouncil,
\emph{A certified uniform descent for all Dirichlet characters},
public manuscript and proof development, 9 October 2026.
Version used: \texttt{40844a7614b67840801e65b836b3da2e1f7ff029}.
\url{https://github.com/GTimG/QuasiRiemannTracker/blob/40844a7614b67840801e65b836b3da2e1f7ff029/public/proofs/qrh-20261009/accepted-answer.tex}.
\end{thebibliography}
\end{document}
